% The script: every word the reader sees, in the order it plays.

% Grammar: GRAMMAR.md. Build the PDF: npm run script:watch (or script:pdf).

\section{Openning}\label{act:openning}

\subsection{Title}\label{title}

The teletype types the news line. Then the title builds in reading order:
MERIT fades in, plain; then "or"; then MATH arrives through its reel; then
"?". Blue sits under MERIT, Red under MATH, breathing. The reel starts on a
couple of funny wrong words, turns slowly enough to read every word, runs
past the answer and comes back to rest on it.

% As today; MERIT has no reel. On a phone the words shrink; they are never
% stacked. In Farsi the whole line mirrors like everything else: MERIT comes
% first, on the right, and Blue with it.

Ledger (teletype): The world has its first trillionaire.
\todo{draft}
\marginpar{The dated source stay attached to the headline.}

\begin{description}
  \item[open.source] (Reuters · June 14, 2026)
\end{description}

% the teletype first, then the crowd hops in, then the title
\crowd{idle}

\reveal{merit}

\reveal{or}

\reveal{math}
\begin{enumerate}
  \item[hold 3] Grit
  \item Stars
  \item Socks
  \item Timing
  \item Friends
  \item Parents
  \item \textbf{Math}
  \item Hair
  \item Blue Eyes
\end{enumerate}
\adapt{Translate each word as the everyday explanation, not literally; keep the answer (in bold) and the words after it wrong.}

% the reel lands on the word in bold (MATH); the words after it are wrong ones
% it overshoots onto before settling back. [hold n] keeps a word n beats. Short words only (owner,
% 2026-09-26): the title is sized so the widest one never leaves the screen,
% so one long word makes the whole title smaller. Race, Genes, God's will and
% Class stay off it: they take political sides.

\subsection{The crowd}\label{crowd}

People come in from different directions, bouncing and jumping, and settle
as a loose crowd — truly random, no rows, no clusters of one shape or
colour, everyone the same size. Coins fall at random: some people catch them
and grow, some do not. Then they bump into each other, and every bump is a
trade. One grows very big, one ends up small. The rest bounce and jump and
move out. Those two stay: the big one under MERIT, the small one under MATH.
No words, no numbers.

\crowd{payout}
\reveal{mark}
\marginpar{Every trade follows the game's own rule — the stake is half of
the smaller fortune, the winner takes it. It ends at 15 and 1.}

\begin{description}
  \item[open.title] Merit or Math?
  \item[open.title.merit] Merit
  \item[open.title.or] or
  \item[open.title.mark] ?
  \item[open.credit] Created \& directed by Hamed
\end{description}

% the common sense comes first, the question second (owner, 2026-09-26).

\adapt{[open.title.mark] is the language's own question mark ("؟" in Farsi).}

\crowd{two}

\section{Introduction}\label{act:intor}

\subsection{Meeting them}\label{meet}

The two survivors, still neutral. Red calls first. A click makes a circle
take its colours and introduce itself at once; then the other one calls. If
the reader waits, the caller tries again, a little louder.

% Both call at once, each at its own random moments; the reader clicks
% either first. Each one's introduction answers the click.
Red (aside, brief):
\begin{enumerate}
  \item Hi.
  \item Hello?
  \item Anybody there?
  \item Click on me.
\end{enumerate}
\hide{headline}
\meet{red}
\todo{draft}

Blue (aside, brief):
\begin{enumerate}
  \item Psst. Over here.
  \item Over here!
  \item Click me. I'm the important one.
\end{enumerate}
\meet{blue}
\todo{draft}

Red (aside, flow): I'm the poorest man in this world.
\on{met-red}

Blue (aside, flow): I'm the richest man in this world.
\on{met-blue}
\marginpar{Blue is LIKE a tech billionaire and is never a real person; he quotes and imitates no one. The essay explains no real person's fortune.}

Red (aside, brief): Click, press a key, or scroll to go on.
\todo{draft}

% the one line about controls, and it waits for the reader. Removed
% 2026-09-26: "'This world' means the one on your screen" and Blue's "Small
% world. But it's mine." ("Everybody knows that.")

\subsection{The merit debate}\label{merit}

Blue argues from under MERIT, Red answers from under MATH.

% Every line here is a key line: it waits for the reader. Red hints at
% chance without a technical word — no coins, no games, no probability.

Blue: I work hard for it.

Red: I work two shifts to pay the bills.

Blue: I'm smart.

Red: I'm smart too.
I have a PhD in physics.

Blue: I take big risks, and win big.

Red: I take risks too.
But nobody catches me when I fall.
\marginpar{These are two characters' opinions, not the essay's findings. Nothing here is shown true until the game shows it.}

% cut 2026-09-26: the GPS pair (too long), "Look, anyone could be where I am"
% with Red's "standing here first…" (out of character; the pair goes), and
% "Success finds the people who deserve it" with its reply.

\subsection{The invitation}\label{invite}

Numbers appear for the first time — each fortune shown as its coins, packed
tight. Every coin is the same size everywhere, and a circle's area is its
coins' area.

Red: Let's play a game.
\reveal{coins}
\todo{draft}

Blue (flow): I don't play games.
I build things.
\todo{draft}

Red: Just try it. We start equal.
\todo{draft}

Blue (flow): Equal?
Why would I give up my wealth?
\todo{draft}

Red: It's only a simulation.

Blue (flow): Fine.
I like a good experiment.

The title and every other word move up and out. Only the two remain.

\hide{words}

Blue (aside): Give him some of my coins.
\equalize
\expect{blue=8, red=8}

A tap on a circle sends one of its coins, visibly, to the other; or drag a
coin across and it follows the finger. Moving coins back is allowed. Each
circle re-sizes by area after every move.

Blue (flow): Hey!
\on{first-move}
\todo{draft}

Blue (flow): I liked that one.
\on{blue-reaches-eleven}
\todo{draft}

Blue (flow): Whoa. Equal, not upside down.
\on{red-above-eight}
\todo{draft}

Red (flow): Tempting. But give him one back.
\on{red-above-eight}
\todo{draft}

Red (flow): Eight and eight. Perfect.
\todo{draft}

\subsection{Round one}\label{round1}

The stakes are the same coins, on the table below the pair. The decider is
one plain coin until it is tossed: Marx is Red's side, the bank is Blue's.

% Blue is the brief's trader A, Red is trader B. The rule is deliberately
% incomplete here.

Red: The rules are simple.
We each put half on the table.

Blue: Half of 8.
4 each.
\stake{4}

Four coins from each circle move to the table.

Red: We flip a coin.
Blue side, you win.
Red side, I win.
\reveal{coin}

Red: The winner takes it all.
\adapt{[r1.winner] echoes a well-known song title in English. The line must still read plainly — the winner of the toss takes both stakes.}

The coin turns; its faces change only when it is edge-on. Red lands. Blue 4,
Red 12.

\flip{red}
\expect{blue=4, red=12}

\begin{description}
  \item[log.toss] \val{winner} wins. Blue \val{blue}, Red \val{red}.
\end{description}

% not a bubble: after every toss the result is logged in the talk, beside the
% coin that landed, and moves up with the lines — so the reader can always see
% what happened (owner, 2026-09-26).

Blue (aside, flow): Ouch!

\subsection{Round two}\label{round2}

Blue: Wait!
I have 4. I can't put in as much as you.
\hide{coin}

Red: Put in half of what you have.

Blue (flow): Only 2?

Red: Yes. And I'll match that.
\stake{2}

Two coins from each circle move to the table.

Blue: So we always put in half
of what the poorer one has.

Red: Yes. That way everybody can keep playing.

Blue (flow): Smart.
Okay, let's go.
I'll win it back.
\reveal{coin}

Blue lands. Blue 6, Red 10.

% Blue says the rule himself and Red confirms it; the stakes are already
% known to be equal, so nobody says "half of each".

\flip{blue}
\expect{blue=6, red=10}

\subsection{Round three}\label{round3}

% slower than before, and the two say what is happening.

Red (flow): 3 each.
\stake{3}
\reveal{coin}
\todo{draft}

Three coins from each circle move to the table.

Blue (brief, flow): Flip!
\todo{draft}

Blue lands. Blue 9, Red 7.

\flip{blue}
\expect{blue=9, red=7}

Blue (aside, flow): And that's how it's done.
\todo{draft}
\marginpar{The three rounds are the real rule played with authored tosses: 8-8, then 4-12, then 6-10, then 9-7. The total never changes, and every stake is half of the poorer fortune.}

\subsection{The challenge}\label{dare}

Blue (aside, flow): Back where I started.
Well, I'm a bit ahead. Naturally.
\hide{coin}

Blue (flow): This is pointless.

Red (flow): Why?
\todo{draft}

Blue: Nobody has an edge.
Same odds, same rule.
Nobody gets rich.

Red: What if one does?
{\Large Very riiich.}

Blue (flow): {\Large Impossible.}

% `**…**` in a bubble is said louder: bigger and heavier.

Red: But what if that happens?

Blue: That being rich means nothing.
No talent, no hard work. Just luck.

Red: Not merit. Just math.
\adapt{[dare.math] must use the SAME two words as the title in each language.}

Blue (flow): Nice line. It won't happen.
\todo{draft}

Red: I can prove it.

Blue: I'll bet half my wealth you can't.

Red: I'll bet all of mine.

Blue (flow): You're either very sure
or very foolish.

Red (flow): Let's see who laughs last.
\todo{draft}

Red: We need more players.
And many more rounds.
\marginpar{A run shows something; it cannot prove it. Red's "prove" is honest only because the proof arrives later — the reveal says so, and the theorem pays it off.}

\section{More players}\label{act:more}

\subsection{The crowd arrives}\label{room}

The room fills around the two: people arrive from all sides and the view
zooms out, the pair keeping its place. The crowd keeps its costumes; Blue
and Red keep their unique colours.

Red: Different colours. Different shapes.
Only the size counts. The size is the wealth.
\crowd{room}
\todo{draft}

Red: Each time,
two people are picked at random.
\marginpar{[more.random] is REQUIRED — the two-person game never taught random pairing, and every room after this one uses it.}

Red: They bet a tenth of what the poorer one has.
Flip. The winner takes it all.
\card{rule}
\pairs{2}
\marginpar{Two rounds played by hand: two from the room, drawn at random, step out, play, and go back. Nothing they do is kept; the run starts the room fresh.}

% 2026-09-27: "a quarter" → "a tenth" — the room's stake is 10% now ("still
% jittery"). 2026-09-26: "half" → "a quarter" — it was 25% (owner: "0.5
% is too large and jittery, maybe .25?"). The two-player game above still
% plays at half.

Blue (flow): Nobody gets rich.

% Blue laughs.

\subsection{The joke offer}\label{offer}

Blue: Real life is
much more complicated than this.

Red (aside): Of course it is.
Want to hear a joke?
\choice{Yes}{joke=yes}
\choice{Not now}{joke=no}

\subsection{The joke}\label{joke}

% Owner, 2026-10-09: the joke is part of the chat. Red tells it in his own
% words and posts the plates with them, like photos in a chat. "Not now"
% skips to the guess. The cow's scroll scene is retired
% (archive/scenes-pre-chat-joke/).

Red (flow, amused): A farmer's cow stopped giving milk.
\picture{introduction}
\todo{draft}

Red (flow): He called a biologist, a chemist and a physicist.
\picture{darwin}
\picture{chemist}
\todo{draft}

Red (flow): They looked at the cow. The cow looked at them.
\picture{silence}
\picture{cow-looks}
\todo{draft}

Red (flow): Moo.
\picture{cow-moos}
\picture{scratch}
\todo{draft}

Red (flow): The physicist: "Assume a spherical cow."
\picture{physicist}
\picture{spherical}
\todo{draft}

Red (flow): "In a vacuum, and no gravity!"
\picture{vacuum}
\todo{draft}

Red (flow): The sphere alone answered this question.
Not every question. This one.
\picture{football}
\todo{draft}

Blue (flow, amused): That's you. That's exactly you.
\todo{draft}

Red (flow, calm): And this game is my spherical cow.
Simple on purpose.
\todo{draft}

Red: Want to see how a person becomes a circle?
\choice{Show me}{human}
\choice{Skip}
\todo{draft}

\section{Your guess}\label{act:guess}

\subsection{Your guess}\label{guess}

The same room. The question, the rule and the choices are all in front of
the reader at once. The rule card opens beside Red's bubble; the choices sit
inside a large bubble.

% No jump to another page (brief Scene 12).

Red (aside): Now it's your turn to guess.
\todo{draft}

Red (aside): Remember the \gls{rule}.
\reveal{card:rule}
\todo{draft}

Red (aside): What do you think will happen?
\choice{Still roughly equal}{prediction=equal}
\choice{A gentle spread}{prediction=spread}
\choice{A split: some rich, some poor}{prediction=split}
\choice{One giant, the rest near nothing}{prediction=giant}
\todo{draft}

\begin{description}
  \item[log.guess] Your guess: \val{choice}.
\end{description}

Red (aside): And how much would you bet on that?
\hide{card:rule}
\choice{A coffee}{bet=coffee}
\choice{A lunch}{bet=lunch}
\choice{A vacation}{bet=vacation}
\choice{1\% of my wealth}{bet=percent}
\todo{draft}

\begin{description}
  \item[log.bet] Your bet: \val{bet}.
\end{description}

Blue (aside, flow): A coffee. Careful with your money. I respect that.
\when{bet=coffee}
\todo{draft}

Red (aside, flow): A lunch. Fair. The winner picks the place.
\when{bet=lunch}
\todo{draft}

Blue (aside, flow): A vacation! Now it's serious.
\when{bet=vacation}
\todo{draft}

Red (aside, flow): One percent. Careful. That's how it starts.
\when{bet=percent}
\todo{draft}

\section{The run}\label{act:run}

\subsection{The run}\label{run}

The same room, the same hundred, Blue and Red among them. One live run,
slowed down so people can see it. The biggest fortune wears the dashed ring.
A forward press finishes the run at once. The result and the morning paper
are logged in the talk.

% Unseeded; slowed down because "it was very fast" (brief).

Red (aside): Here we go.
Fair trades, one after another.
\todo{draft}

\run

\begin{description}
  \item[log.run] After \val{trades} trades, one of them holds \val{share} of everything.
\end{description}

% after the run, in the talk; then the morning paper prints there too, on the
% winner, as today.

Blue (flow): Red? Where are you?
\when{winner=other}

Red (flow): Here. What do you have?
\when{winner=other}

Blue (flow): \val{blue}. You?
\when{winner=other}
\todo{draft}

Red (flow): \val{red}.
\when{winner=other}
\todo{draft}

Red (flow): Told you.
\when{winner=other}

% when someone else finishes richest (the most likely case); {blue} and {red}
% are what they actually hold, in dollars (everyone started with $100).

Blue (flow): See? Talent rises.
\when{winner=blue}
\todo{draft}

Red (flow): Run it again and see if your talent comes along.
\when{winner=blue}
\todo{draft}

% when Blue finishes richest.

Red (flow): Look at that. Am I a genius now?
\when{winner=red}
\todo{draft}

Blue (flow): …Run it again.
\when{winner=red}
\todo{draft}

% when Red finishes richest.

Red (aside): Let's try again?
\choice{Again}{\run}
\choice{Go on}
\todo{draft}

\subsection{So why did they win?}\label{why}

The paper is in the talk. Blue takes the paper's side; Red refuses.

% The bias is described, never named (owner, 2026-09-26; research.md C12).

Blue: See? The paper knows why.
\todo{draft}

Red: The paper found the reason after the result.
The coin never saw a thing.
\when{runs=one}
\todo{draft}

Red: New winner. New reason.
The paper never prints "chance."
\when{runs=several}
\todo{draft}

% instead of [why.after] once the reader has run the room more than once.

% Red's one line about the reader's last room, whichever fits.

Red (aside): And that's one run. It shows, it doesn't prove.
The proof comes later.
\todo{draft}
\marginpar{Every claim here is about one finite run. No "always", no distribution family, no naming a psychological bias. The result comes first and the explanation second; colour, corners and edges never touch the game.}

% OPEN — the payoff of Blue and Red's bet, and of the reader's.

\section{Line them up}\label{act:sort}

\subsection{Line them up}\label{sort}

The same room. The people drop and sort into piles by how much they hold,
keeping their costumes; the count is written above each pile; the ruler is
ordinary, its ticks round. Then the ruler changes to multiplying by ten, a
touch slower, so the labels can be seen moving. Then everyone goes back.

Red (aside): A hundred people.
Too many to read at a glance.
\choice{Sort them}
\todo{draft}

\arrange{piles}

Blue: Most have almost nothing. A few have a lot.
That's just real life.
\todo{draft}

Red: It's also this room.
And here, nobody had an edge.
\todo{draft}

Blue (flow): Where am I?
\todo{draft}

Red (flow): There. With \val{count} others.
\todo{draft}

% {count} is how many others share Blue's pile, off the screen.

Blue: This chart is useless, though.
Everyone is squeezed into one corner.
\todo{draft}

Red: Then let's change the ruler.
\choice{Change the ruler}
\todo{draft}

\arrange{ruler}

Red: Now each step means ten times more,
not ten more.
\todo{draft}

Blue (flow): One, ten, a hundred, a thousand.
Even steps. Huh.
\todo{draft}

Red: Less than a cent goes in the dust box.
Zero has no place on this ruler.
\todo{draft}

Red (aside): That's a \gls{histogram}.
Back to the room?
\choice{Put them back}
\todo{draft}

\arrange{free}
\card{histogram}
\pin{histogram}

\section{Measure the room}\label{act:gini}

\subsection{Measure the room}\label{gini}

Everyone steps into a line, poorest first. Walking along the line, their
money is added up as it goes: the running total climbs as a curve above
them. Equal shares would climb the straight diagonal. The gap between the
two is the Gini.

Blue: Economists love a single number.
What's ours?
\todo{draft}

Red: Line everyone up, poorest first.
\arrange{line}
\todo{draft}

Red: Now walk along the line
and add up their money as you go.
\reveal{curve}
\todo{draft}

Red: If everyone had the same,
the total would climb this straight line.
\reveal{diagonal}
\todo{draft}

Red: The gap between them is the \gls{gini}{Gini}.
Zero: all the same. One: a single owner.
\reveal{gap}
\todo{draft}

Red: This room: \val{gini}.
\todo{draft}

Red (aside): Want to see the equal case
and play with it?
\choice{Yes}{\reveal{toy:gini}}
\choice{Not now}
\todo{draft}

\arrange{free}
\hide{lorenz}
\card{gini}
\pin{gini}

\section{How many still count?}\label{act:count}

\subsection{How many still count?}\label{count}

The room itself tries the cases: everyone the same; one person emptied; that
person's money given to one other; half the room owning it all; one owner;
and back to how it really is. Each time Red says the number. Then four
people step out of the crowd, four coins each, and the reader moves the
coins.

% Owner, 2026-09-26.

Blue: A hundred players.
That's a lot of competition.
\todo{draft}

Red: If all hundred had the same,
all hundred would count.
\arrange{equal}
\todo{draft}

Red: This game is brutal.
No money, and you don't count at all. \val{count} left.
\arrange{zero}
\todo{draft}

Red: Now give his money to one other person.
\val{count}. A bit less.
\arrange{double}
\todo{draft}

Red: Half of them own it all, equally?
\val{count}.
\arrange{half}
\todo{draft}

Red: One owns everything?
\val{count}.
\arrange{one}
\todo{draft}

Red: This room?
About \val{count}.
\arrange{free}
\todo{draft}

% {count} is 1 / Σsᵢ² of the room on screen, to one decimal.

Red (aside): Your turn. Four of them, four coins each.
Move the coins.
\arrange{four}
\choice{Done}
\todo{draft}

Red (aside, interrupts): Four people. \val{count} of them count.
\todo{draft}

% Red, beside the four, reading the number out as the reader moves coins.

Blue: A hundred people.
About \val{count} players.
\arrange{free}
\card{participants}
\pin{participants}
\todo{draft}

\section{Is anything moving?}\label{act:turnover}

\subsection{Is anything moving?}\label{turnover}

The room dims; over it, how much of the room's money changed hands in each
round of the run, from the first to the last.

% Owner, 2026-09-26: "people want an active economy; low turnover is bad
% even for the right".

Blue: At least it's a busy economy.
\val{trades} trades!
\todo{draft}

Red: Count the money that actually changes hands.
\arrange{turnover}
\todo{draft}

Red: At first, \val{early} of the room changed hands every round.
\todo{draft}

Red: By the end, \val{late}.
The trades go on. The money barely moves.
\todo{draft}

Blue: A dead economy.
Even I don't like that.
\todo{draft}

\arrange{free}
\card{turnover}
\pin{turnover}

\section{Where it ends}\label{act:end}

\subsection{Where it ends}\label{end}

The same room. Red states what the mathematics says; "Run it longer?" plays
the same room on, fast, in place: closer and closer to one owner, never on a
date.

% The 1,000-person crowd is cut from the main flow (D20 answer 8).

Red (aside): That was one run.
It shows. It doesn't prove.
\todo{draft}

Red: The math says more.
Keep playing forever, and one person ends up with everything.
\todo{draft}

Blue (flow): When?
\todo{draft}

Red: Never on a date. It's a \gls{limit}.
At any moment, it can still wobble.
\todo{draft}

Red (aside): Want to see it run longer?
\card{limit}
\choice{Yes}{\run[longer]}
\choice{Not now}
\todo{draft}

\section{Your hand on the dial}\label{act:dial}

\subsection{Your hand on the dial}\label{dial}

A stake dial inside Red's bubble, from nothing to everything. Every change
runs a fresh room of the same length, so luck is allowed to heckle the
comparison.

% Brief 1.5: nothing is capped.

Red (aside): Your turn.
Pick a stake, from nothing to everything.
\control{stake}
\todo{draft}

Red (flow): Zero: nothing moves.
That's no game at all.
\when{stake=zero}
\todo{draft}

Red (flow): \val{stake}: slower.
The richest holds \val{share}.
\when{stake=slow}
\todo{draft}

Red (flow): \val{stake}: the stake you watched.
The richest holds \val{share}.
\when{stake=same}
\todo{draft}

Red (flow): \val{stake}: faster.
The richest holds \val{share}.
\when{stake=fast}
\todo{draft}

Red (flow): Everything on the table:
one loss and you're out. The richest holds \val{share}.
\when{stake=all}
\todo{draft}

Red: Every \gls{stake} above zero ends the same way.
The stake changes the journey, not the end.
\control{none}
\card{stake}
\todo{draft}
\marginpar{The qualitative split only — every fixed stake in (0, 1) has the same limit (C4); the rate is what these runs did, not a law. Exactly zero switches the game off.}

\section{Now you try to stop it}\label{act:stop}

\subsection{Now you try to stop it}\label{stop}

The room trades live. Tapping a fortune takes a quarter of it into a pool
and shares it back equally — a manual WEALTH levy. Keep at least twenty
effective participants for thirty seconds.

% C13. The pace is tuned so a diligent hand genuinely can.

Red (aside): Enough watching.
Keep at least twenty players in the game.
\control{tax}
\todo{draft}

Red: Tap a big fortune: a quarter goes into a pool
and comes back to everyone, equally.
\choice{Start}{\run[game]}
\todo{draft}

Red: Thirty seconds, and still \val{count} players.
Your hand did it.
\on{game-won}
\todo{draft}

Red: It closed after \val{seconds} seconds.
\val{taps} taps weren't enough.
\on{game-lost}
\todo{draft}

Red: Notice what your hand did:
watch, pick one, react. Again and again.
\control{none}
\todo{draft}

Red: Next, not a hand. A rule.
Nobody gets picked.
\todo{draft}

\section{Put the levy in the rules}\label{act:levy}

\subsection{Put the levy in the rules}\label{levy}

Trading frozen; four step out with coins — Blue 16, two others 8 and 4, Red
4 (the average is 8). A quarter of every pile goes to one pool; the pool
comes back in equal parts: Blue 14, 8, 5, Red 5.

Red (aside): Your hand gets tired.
A rule doesn't.
\arrange{levy4}
\todo{draft}

Red: Once a round, the same share from everyone:
a quarter.
\levy{collect}
\todo{draft}

Blue (flow): Everyone lost a quarter.
Nothing changed.
\todo{draft}

Red: Right. Everyone shrank the same.
The shares didn't move.
\todo{draft}

Red: Now the pool goes back,
in equal parts.
\levy{return}
\todo{draft}
\marginpar{The return is literal coins here; outside the room it could be a dividend or a shared service — the model chooses no budget, it only tests the shared return.}

Red: The average broke even.
Below it, more came back. Above it, less.
\todo{draft}

Blue (aside, flow): I paid in.
\todo{draft}

Red: And the room kept every coin.
\card{levy}
\todo{draft}

\section{Trade and return together}\label{act:match}

\subsection{Trade and return together}\label{match}

The room shrinks to one side and a mirror copy appears beside it — a
parallel universe with the same start and the same luck. Left: trades only.
Right: the same trades plus a 0.3\% levy and equal return every round.

% 375,000 trades at a tenth: about 4 still count on the left, about
% 41 on the right.

Red: Now let it trade. Same luck, twice:
one room without the rule, one with it.
\match
\todo{draft}

Red: Without the rule: about \val{a} players.
With it: about \val{b}.
\todo{draft}

Blue (flow): Same coin, same partners?
\todo{draft}

Red: Every toss the same.
Only the rule changed.
\todo{draft}
\marginpar{One matched pair is one finite experiment; within the pair luck is held still, and the only changed ingredient is the levy plus equal return.}

\section{The outcome map}\label{act:map}

\subsection{The outcome map}\label{map}

The map of stake against levy, each square a finished room, coloured by how
many still count. The dashed fit is drawn only over the range it was
measured on; neither it nor the contour is a phase boundary.

Blue: One pair is confusing.
Which settings keep the room open?
\arrange{free}
\todo{draft}

Red: Try them all: stake across, levy up.
Each square, how many still count.
\reveal{map}
\todo{draft}

Red: A small levy can hold a big stake:
double the stake, about four times the levy.
\reveal{fit}
\todo{draft}
\marginpar{The fit is descriptive, drawn only over the measured range; each square is the mean of 8 finished rooms (200,000 trades, levy once a round); the stake is at risk each trade, the levy comes once a round — different clocks, compared by their measured effects, not interchangeable percentages.}

\section{The verdict}\label{act:verdict}

\subsection{The verdict}\label{verdict}

Everyone else hops out of the room; Blue and Red go back to their seats,
eight coins each, as the reader once made them — equal.

% Brief 5.10: "everything leaves except the two. They talk".

Red: The coin was fair the whole time. It never saw a name.
\crowd{empty}
\expect{blue=8, red=8}
\todo{draft}

Blue: Then how did one of us end up with everything?
\todo{draft}

Red: Money didn't buy better odds. It bought more tosses you could survive.
\todo{draft}

Red: Then a shared levy with an equal return kept more people able to play.
\todo{draft}

Blue (flow): So you're against merit.
\todo{draft}

Red: No. I'm worried about the field merit plays on.
\todo{draft}

Blue: If merit is supposed to win…
\todo{draft}

Red: …keep the game open long enough for merit to play.
\todo{draft}
\marginpar{The conclusion distinguishes fair odds from staying power, names the full levy-plus-equal-return mechanism, explains no real person's wealth, and is not an argument against merit.}

% proposed: the essay's closing line goes to the merit defender.

% OPEN — the ending as a whole, including the payoff of both bets.

\section{The machine is yours}\label{act:sandbox}

\subsection{The machine is yours}\label{sandbox}

The hundred hop back in, equal, Blue and Red among them. The room is the
reader's now: Play, the stake, the levy and how often, the speed, and what a
tap does, in the right panel beside the plots; a tap on a fortune levies it
or photographs it for the paper. "All the dials" opens the whole old machine
(people, money, expert numbers, the measured map) as a side trip.

% Brief 5.10: "Blue and Red stay in the sandbox too".

Red: I'm done touching the controls.
\crowd{room}
\control{sandbox}
\todo{draft}

Blue (aside): The machine is yours. Break our argument.
\choice{All the dials}{workshop}
\todo{draft}
